\documentclass[11pt,twocolumn]{article}

% === PACKAGES ===
\usepackage[utf8]{inputenc}
\usepackage[T1]{fontenc}
\usepackage{mathpazo}
\usepackage{amsmath,amssymb,amsthm}
\usepackage{graphicx}
\usepackage{xcolor}
\usepackage{hyperref}
\usepackage{cleveref}
\usepackage{booktabs}
\usepackage[margin=0.75in]{geometry}
\usepackage{enumitem}
\usepackage{multirow}
\usepackage{longtable}

% === COLORS ===
\definecolor{linkblue}{RGB}{0,51,153}

% === HYPERREF SETUP ===
\hypersetup{
    colorlinks=true,
    linkcolor=linkblue,
    citecolor=linkblue,
    urlcolor=linkblue
}

% === NOTATION ===
\newcommand{\eps}{\varepsilon}
\newcommand{\eM}{$\eps$-machine}
\newcommand{\Cmu}{C_\mu}
\newcommand{\hmu}{h_\mu}
\newcommand{\TargetRefs}{120}

\pdfstringdefDisableCommands{%
    \def\eps{epsilon}%
    \def\eM{epsilon-machine}%
    \def\Cmu{Cmu}%
    \def\hmu{hmu}%
}

% === TITLE ===
\title{%
    \vspace{-1em}
    {\Large Computational Mechanics: A State-of-the-Art Review (1989--2026)}\\[0.3em]
    {\normalsize Theory, Inference, Extensions, and Applications}
}
\author{%
    John S Azariah\\
    \small University of Technology, Sydney (UTS)\\
    \small\texttt{john.azariah@student.uts.edu.au}
}
\date{February 2026}

\begin{document}

\twocolumn[
\maketitle

\begin{abstract}
This manuscript is a clean-slate, comprehensive review of Computational Mechanics designed as a field checkpoint rather than a tutorial.
The review synthesises the evolution of foundational theory, inference algorithms, formal extensions, and domain applications from 1989 to 2026.
Its central objective is to map how the field has changed over time, what has been validated, where methods are now used, and which open problems remain unresolved.
The final version will include at least \TargetRefs{} curated references with explicit coverage across foundations, methods, extensions, applications, and recent advances.
\end{abstract}
\vspace{1.2em}
]

% ===================================================================
\section{Scope and Review Protocol}
\label{sec:protocol}
% ===================================================================

\subsection{Review Questions}

This review is organised around four questions.
First, how has the core formalism of causal states and \eM s evolved since the original formulation.
Second, which inference methods are now considered mature, and under what data regimes.
Third, which extensions beyond the classical finite-alphabet stationary setting are established.
Fourth, what evidence exists for real-world utility across scientific domains.

\subsection{Inclusion Criteria and Evidence Quality}

Included literature must contribute at least one of the following: foundational theory, methodological innovation, rigorous benchmarking, or substantive domain application.
Each section will distinguish between conceptual proposals, mathematically proven results, and empirically validated outcomes.

\subsection{Coverage Targets}

The reference budget is distributed across categories to avoid over-concentration on foundational papers.

\begin{table}[h]
\centering
\caption{Minimum reference coverage targets}
\label{tab:ref-targets}
\small
\begin{tabular}{lcc}
\toprule
Category & Minimum & Status \\
\midrule
Foundations and core theory & 15 & TODO \\
Inference algorithms and estimation & 25 & TODO \\
Formal extensions (quantum, multivariate, continuous) & 25 & TODO \\
Applications across domains & 35 & TODO \\
Recent results (2020--2026) & 20 & TODO \\
\midrule
Total unique references & $\geq \TargetRefs$ & TODO \\
\bottomrule
\end{tabular}
\end{table}

% ===================================================================
\section{Foundational Trajectory of the Field}
\label{sec:foundations}
% ===================================================================

This section will present a chronological synthesis of the formal core.
The emphasis is on identifying which theorems and definitions remained stable, which were refined, and how interpretive framing changed across decades.
Rather than reproducing long proofs, the review will extract the enduring results and their modern implications.

\paragraph{Planned synthesis blocks.}
\begin{itemize}[nosep]
    \item Origins of causal-state reconstruction and predictive sufficiency.
    \item Canonical results on minimality, unifilarity, and uniqueness.
    \item Clarifications and reinterpretations in later literature.
    \item Relationship between $E$, $\Cmu$, crypticity, and predictive information.
\end{itemize}

% ===================================================================
\section{Inference Algorithms: Evolution and Trade-offs}
\label{sec:inference}
% ===================================================================

This section will compare algorithmic families as they developed over time.
Each family will be reviewed by assumptions, computational complexity, finite-sample behaviour, robustness, and practical tuning burden.

\subsection{Algorithm Families}

\begin{itemize}[nosep]
    \item CSSR and descendants.
    \item Spectral and operator-based approaches.
    \item Bayesian structural inference methods.
    \item Neural and representation-learning based state discovery.
\end{itemize}

\begin{table*}[t]
\centering
\caption{Algorithm comparison template for literature synthesis}
\label{tab:algo-synthesis}
\small
\begin{tabular}{lllllll}
\toprule
Family & Typical assumptions & Data regime & Runtime scaling & Strengths & Failure modes & Key papers \\
\midrule
CSSR-like & TODO & TODO & TODO & TODO & TODO & TODO \\
Spectral & TODO & TODO & TODO & TODO & TODO & TODO \\
BSI-like & TODO & TODO & TODO & TODO & TODO & TODO \\
Neural/latent & TODO & TODO & TODO & TODO & TODO & TODO \\
\bottomrule
\end{tabular}
\end{table*}

% ===================================================================
\section{Formal Extensions Beyond the Classical Setting}
\label{sec:extensions}
% ===================================================================

This section will review how Computational Mechanics has expanded beyond finite, stationary, univariate symbolic processes.

\paragraph{Extension axes.}
\begin{itemize}[nosep]
    \item Quantum and q-machine formulations.
    \item Continuous-valued and mixed observation settings.
    \item Spatiotemporal and multivariate generalisations.
    \item Non-stationary and adaptive formulations.
\end{itemize}

For each axis, the review will explicitly mark maturity level: theoretical proposal, partial validation, or well-established practice.

% ===================================================================
\section{Application Landscape and Cross-domain Uptake}
\label{sec:applications}
% ===================================================================

A review paper should map where the framework has actually been used, not only where it could be used.
This section therefore synthesises applications by domain and evaluates evidence quality.

\begin{table*}[t]
\centering
\caption{Application landscape template}
\label{tab:application-landscape}
\small
\begin{tabular}{llllll}
\toprule
Domain & Typical data & CM objective & Representative studies & Evidence quality & Open gap \\
\midrule
Physics & TODO & TODO & TODO & TODO & TODO \\
Neuroscience & TODO & TODO & TODO & TODO & TODO \\
Biology/Genomics & TODO & TODO & TODO & TODO & TODO \\
Linguistics & TODO & TODO & TODO & TODO & TODO \\
Finance/Economics & TODO & TODO & TODO & TODO & TODO \\
Engineering/Control & TODO & TODO & TODO & TODO & TODO \\
\bottomrule
\end{tabular}
\end{table*}

% ===================================================================
\section{What Changed in 2020--2026}
\label{sec:recent}
% ===================================================================

This section will isolate genuinely new results from incremental re-packaging.
The goal is to state clearly what is now known that was not established in earlier decades.

\paragraph{Planned outputs.}
\begin{itemize}[nosep]
    \item Timeline of major post-2020 contributions.
    \item Which claims received independent replication.
    \item Which proposed methods remain unvalidated.
    \item Which application domains saw first credible uptake.
\end{itemize}

% ===================================================================
\section{Open Problems and Research Agenda}
\label{sec:agenda}
% ===================================================================

This section will convert the synthesis into an actionable research map.
Open problems will be prioritised by impact, tractability, and required prerequisites.

\begin{table*}[t]
\centering
\caption{Open-problem matrix template}
\label{tab:open-problems}
\small
\begin{tabular}{llllll}
\toprule
Problem & Why unsolved & Best current approach & Evidence status & Blocker & Near-term experiment \\
\midrule
Finite-sample uncertainty quantification & TODO & TODO & TODO & TODO & TODO \\
Non-stationary inference & TODO & TODO & TODO & TODO & TODO \\
Scalable multivariate reconstruction & TODO & TODO & TODO & TODO & TODO \\
Quantum advantage under noise & TODO & TODO & TODO & TODO & TODO \\
\bottomrule
\end{tabular}
\end{table*}

% ===================================================================
\section{Conclusion}
\label{sec:conclusion}
% ===================================================================

The final review will function as a checkpoint of record for the field.
Its deliverable is not only a summary of concepts, but a structured map of how Computational Mechanics evolved, where it is now applied, and what evidence supports current claims.
Completing the literature synthesis to at least \TargetRefs{} references is therefore a core methodological requirement rather than a cosmetic target.

% ===================================================================
% ACKNOWLEDGMENTS
% ===================================================================
\section*{Acknowledgments}

GitHub Copilot (GPT-5.3-Codex and Claude) was used during development for
code generation, refactoring, test scaffolding, and documentation drafting.
All AI-generated outputs were reviewed, edited, and validated by the author,
who made all core research and design decisions.


% ===================================================================
% BIBLIOGRAPHY
% ===================================================================
\bibliographystyle{plain}
\bibliography{../shared/bibliography/references}

\end{document}
